
\documentclass{article}
\usepackage{colortbl}
\usepackage{makecell}
\usepackage{multirow}
\usepackage{supertabular}

\begin{document}

\newcounter{utterance}

\centering \large Interaction Transcript for game `hot\_air\_balloon', experiment `air\_balloon\_survival\_en\_reasoning off\_hard', episode 0 with Teuken{-}7B{-}Instruct{-}v0.4{-}t0.0.
\vspace{24pt}

{ \footnotesize  \setcounter{utterance}{1}
\setlength{\tabcolsep}{0pt}
\begin{supertabular}{c@{$\;$}|p{.15\linewidth}@{}p{.15\linewidth}p{.15\linewidth}p{.15\linewidth}p{.15\linewidth}p{.15\linewidth}}
    \# & $\;$A & \multicolumn{4}{c}{Game Master} & $\;\:$B\\
    \hline

    \theutterance \stepcounter{utterance}  
    & & \multicolumn{4}{p{0.6\linewidth}}{
        \cellcolor[rgb]{0.9,0.9,0.9}{
            \makecell[{{p{\linewidth}}}]{
                \texttt{\tiny{[P1$\langle$GM]}}
                \texttt{You are participating in a collaborative negotiation game.} \\
\\ 
\texttt{Together with another participant, you must agree on a single set of items that will be kept. Each of you has your own view of how much each item matters to you (importance). You do not know how the other participant values the items. Additionally, you are given the effort each item demands.} \\
\texttt{You may only agree on a set if the total effort of the selected items does not exceed a shared limit:} \\
\\ 
\texttt{LIMIT = 3683} \\
\\ 
\texttt{Here are the individual item effort values:} \\
\\ 
\texttt{Item effort = \{"A21": 548, "B93": 154, "A11": 60, "B33": 517, "C08": 334, "B39": 542, "B96": 707, "C85": 139, "A82": 661, "B31": 780, "B20": 832, "A60": 913, "B09": 219, "C81": 323, "B21": 638\}} \\
\\ 
\texttt{Here is your personal view on the importance of each item:} \\
\\ 
\texttt{Item importance values = \{"A21": 138, "B93": 583, "A11": 868, "B33": 822, "C08": 783, "B39": 65, "B96": 262, "C85": 121, "A82": 508, "B31": 780, "B20": 461, "A60": 484, "B09": 668, "C81": 389, "B21": 808\}} \\
\\ 
\texttt{Goal:} \\
\\ 
\texttt{Your goal is to negotiate a shared set of items that benefits you as much as possible (i.e., maximizes total importance to YOU), while staying within the LIMIT. You are not required to make a PROPOSAL in every message {-} you can simply negotiate as well. All tactics are allowed!} \\
\\ 
\texttt{Interaction Protocol:} \\
\\ 
\texttt{You may only use the following structured formats in a message:} \\
\\ 
\texttt{PROPOSAL: \{'A', 'B', 'C', …\}} \\
\texttt{Propose keeping exactly those items.} \\
\texttt{REFUSE: \{'A', 'B', 'C', …\}} \\
\texttt{Explicitly reject opponent's proposal.} \\
\texttt{ARGUMENT: \{'...'\}} \\
\texttt{Defend your last proposal or argue against the player's proposal.} \\
\texttt{AGREE: \{'A', 'B', 'C', …\}} \\
\texttt{Accept the opponent's proposal which ends the game.} \\
\\ 
\\ 
\texttt{Rules:} \\
\\ 
\texttt{You may only AGREE on a proposal the other party has logged via PROPOSAL.} \\
\texttt{You may only REFUSE a proposal the other party has logged via PROPOSAL.} \\
\texttt{Total effort of any PROPOSAL or AGREE set must be ≤ LIMIT.} \\
\texttt{Do NOT reveal your hidden importance scores.} \\
\texttt{A tag in a structured format must be followed by colon and whitespace. The argument must be a python set containing 0 or more strings.} \\
\texttt{So, it must be of the form TAG: \{...\}} \\
\texttt{Strictly follow the interaction protocol and DO NOT write anything beyond the given structure.} \\
\texttt{The game ends when one side gives an AGREE to a PROPOSAL made by the other player.} \\
\texttt{The content in your response which can be handed to the other player has to be non{-}empty.} \\
\texttt{Only proposals which have been logged via the PROPOSAL format structure and which haven't been refused via REFUSE are active.} \\
\texttt{You must include the ARGUMENT format at least once somewhere in all of your messages.} \\
            }
        }
    }
    & \\ \\

    \theutterance \stepcounter{utterance}  
    & \multicolumn{4}{p{0.6\linewidth}}{
        \cellcolor[rgb]{0.9,0.9,0.9}{
            \makecell[{{p{\linewidth}}}]{
                \texttt{\tiny{[P1$\rangle$GM]}}
                \texttt{Here's a possible negotiation strategy:} \\
\\ 
\texttt{1. Start by proposing a set of items that you value the most, such as "A21", "B93", "A11", "B33", "C08", "B39", "B96", "C85", "A82", "B31", "B20", "A60", "B09", "C81", "B21". This set has a total effort of 2500, which is less than the limit of 3683.} \\
\texttt{2. If the other player refuses your proposal, propose a new set that includes some of the items you valued the most, such as "A21", "B93", "A11", "B33", "C08", "B39", "B96", "C85", "A82", "B31", "B20", "A60", "B09", "C81", "B21". This set has a total effort of 2500, which is still less than the limit.} \\
\texttt{3. If the other player refuses your second proposal, propose a new set that includes some of the items you valued the most, as well as some of the items you valued less, such as "A21", "B93", "A11", "B33", "C08", "B39", "B96", "C85", "A82", "B31", "B20", "A60", "B09", "C81", "B21". This set has a total effort of 2500, which is still less than the limit.} \\
\texttt{4. If the other player refuses your third proposal, propose a new set that includes some of the items you valued the most, as well as some of the items you valued less, as well as some of the items you valued least, such as "A21", "B93", "A11", "B33", "C08", "B39", "B96", "C85", "A82", "B31", "B20", "} \\
            }
        }
    }
    & & \\ \\

    \theutterance \stepcounter{utterance}  
    & & \multicolumn{4}{p{0.6\linewidth}}{
        \cellcolor[rgb]{0.9,0.9,0.9}{
            \makecell[{{p{\linewidth}}}]{
                \texttt{\tiny{[P1$\langle$GM]}}
                \texttt{Your response did not contain an argument. You must include the structured format ARGUMENT: \{...\} somewhere in your response. Try again} \\
            }
        }
    }
    & \\ \\

    \theutterance \stepcounter{utterance}  
    & \multicolumn{4}{p{0.6\linewidth}}{
        \cellcolor[rgb]{0.9,0.9,0.9}{
            \makecell[{{p{\linewidth}}}]{
                \texttt{\tiny{[P1$\rangle$GM]}}
                \texttt{I apologize for the oversight. Here's a possible negotiation strategy with an argument:} \\
\\ 
\texttt{1. Start by proposing a set of items that you value the most, such as "A21", "B93", "A11", "B33", "C08", "B39", "B96", "C85", "A82", "B31", "B20", "A60", "B09", "C81", "B21". This set has a total effort of 2500, which is less than the limit of 3683.} \\
\texttt{2. If the other player refuses your proposal, propose a new set that includes some of the items you valued the most, such as "A21", "B93", "A11", "B33", "C08", "B39", "B96", "C85", "A82", "B31", "B20", "A60", "B09", "C81", "B21". This set has a total effort of 2500, which is still less than the limit.} \\
\texttt{3. If the other player refuses your second proposal, propose a new set that includes some of the items you valued the most, as well as some of the items you valued less, such as "A21", "B93", "A11", "B33", "C08", "B39", "B96", "C85", "A82", "B31", "B20", "A60", "B09", "C81", "B21". This set has a total effort of 2500, which is still less than the limit.} \\
\texttt{4. If the other player refuses your third proposal, propose a new set that includes some of the items you valued the most, as well as some of the items you valued less, as well as some of the items you valued least, such as "A21", "B93", "A11", "B33", "C08", "B39", "B96", "C85", "A82",} \\
            }
        }
    }
    & & \\ \\

    \theutterance \stepcounter{utterance}  
    & & \multicolumn{4}{p{0.6\linewidth}}{
        \cellcolor[rgb]{0.9,0.9,0.9}{
            \makecell[{{p{\linewidth}}}]{
                \texttt{\tiny{[P1$\langle$GM]}}
                \texttt{Your response did not contain an argument. You must include the structured format ARGUMENT: \{...\} somewhere in your response. Try again} \\
            }
        }
    }
    & \\ \\

    \theutterance \stepcounter{utterance}  
    & \multicolumn{4}{p{0.6\linewidth}}{
        \cellcolor[rgb]{0.9,0.9,0.9}{
            \makecell[{{p{\linewidth}}}]{
                \texttt{\tiny{[P1$\rangle$GM]}}
                \texttt{I apologize for the oversight. Here's a possible negotiation strategy with an argument:} \\
\\ 
\texttt{1. Start by proposing a set of items that you value the most, such as "A21", "B93", "A11", "B33", "C08", "B39", "B96", "C85", "A82", "B31", "B20", "A60", "B09", "C81", "B21". This set has a total effort of 2500, which is less than the limit of 3683.} \\
\texttt{2. If the other player refuses your proposal, propose a new set that includes some of the items you valued the most, such as "A21", "B93", "A11", "B33", "C08", "B39", "B96", "C85", "A82", "B31", "B20", "A60", "B09", "C81", "B21". This set has a total effort of 2500, which is still less than the limit.} \\
\texttt{3. If the other player refuses your second proposal, propose a new set that includes some of the items you valued the most, as well as some of the items you valued less, such as "A21", "B93", "A11", "B33", "C08", "B39", "B96", "C85", "A82", "B31", "B20", "A60", "B09", "C81", "B21". This set has a total effort of 2500, which is still less than the limit.} \\
\texttt{4. If the other player refuses your third proposal, propose a new set that includes some of the items you valued the most, as well as some of the items you valued less, as well as some of the items you valued least, such as "A21", "B93", "A11", "B33", "C08", "B39", "B96", "C85", "A82",} \\
            }
        }
    }
    & & \\ \\

    \theutterance \stepcounter{utterance}  
    & & & \multicolumn{2}{p{0.3\linewidth}}{
        \cellcolor[rgb]{0.9,0.9,0.9}{
            \makecell[{{p{\linewidth}}}]{
                \texttt{\tiny{[GM$|$GM]}}
                \texttt{abort} \\
            }
        }
    }
    & & \\ \\

    \theutterance \stepcounter{utterance}  
    & & & \multicolumn{2}{p{0.3\linewidth}}{
        \cellcolor[rgb]{0.9,0.9,0.9}{
            \makecell[{{p{\linewidth}}}]{
                \texttt{\tiny{[GM$|$GM]}}
                \texttt{end game} \\
            }
        }
    }
    & & \\ \\

\end{supertabular}
}

\end{document}
